% ============================================
% MEM-02
% Gerado de molambudos.md — build_miolo.py
% ============================================
\fragtitle{mem}{M}{MEM}{02}{A Caminhada dos Retirantes}
\markboth{MEM-02 --- A Caminhada dos Retirantes}{}

\noindent \textit{1915 --- A estrada de Senador Pompeu a Fortaleza}

\noindent \textbf{1915 --- A estrada de Senador Pompeu a Fortaleza}

\vspace{0.5em}\noindent\rule{\textwidth}{0.5pt}\vspace{0.5em}

\noindent Depois de meses de retirada desde Quixeramobim, chegamos a Senador Pompeu. Em 1915, partimos dali na direção de Fortaleza, seguindo o rumor de que no litoral ainda havia comida.

\noindent No primeiro dia, a gente andou até o sol se pôr.

\noindent Meu pai na frente. Minha mãe atrás, com o saco de roupa na cabeça. Eu no meio, segurando a barra da saia dela. O chão era seco e rachado, e cada passo levantava uma nuvenzinha de poeira que entrava no nariz e deixava a boca com gosto de terra.

\noindent De vez em quando a gente passava por uma casa. As portas estavam abertas. As janelas também. Não tinha ninguém dentro. As casas estavam vazias como cascas de cigarra.

\noindent No segundo dia, a gente começou a encontrar os outros.

\noindent Eles vinham de todo lado. Famílias inteiras. Gente sozinha. Velhos carregando crianças. Crianças carregando velhos. Todo mundo andando na mesma direção.

\noindent Ninguém falava.

\noindent Falar gasta saliva. Saliva precisa de água. Água não tinha.

\noindent A gente andava.

\noindent No terceiro dia, a fila era tão grande que eu não via o fim. A estrada estava cheia de gente que se movia devagar, como um rio de lama. As pessoas tropeçavam umas nas outras. Algumas caíam e não levantavam. A gente passava por cima.

\noindent Minha mãe disse: \textit{\textquotedblleft{}Não olha, Joaquim.\textquotedblright{}}

\noindent Eu olhava.

\noindent No quarto dia, meu pai começou a andar diferente. Os passos ficaram mais curtos. A respiração ficou mais alta. Ele parou de olhar pra frente. Ele olhava para o chão, como se o chão tivesse uma resposta.

\noindent \textit{\textquotedblleft{}Pai\textquotedblright{}\textit{, eu disse. }\textquotedblleft{}Pai, eu quero água.\textquotedblright{}}

\noindent Ele não respondeu.

\noindent \textit{\textquotedblleft{}Pai.\textquotedblright{}}

\noindent Ele parou. Virou o rosto. A boca dele estava aberta. A língua estava branca. Os olhos estavam vermelhos.

\noindent \textit{\textquotedblleft{}Pai?\textquotedblright{}}

\noindent Ele tentou falar. Não saiu som. Só ar. Um ar que cheirava a estômago vazio.

\noindent Ele sentou no chão.

\noindent Minha mãe sentou do lado dele.

\noindent Eu sentei do lado dela.

\noindent A fila de retirantes passou por nós. Ninguém parou. Ninguém perguntou. Todo mundo estava ocupado morrendo.

\noindent No quinto dia, ele parou debaixo de um juazeiro seco.

\noindent Minha mãe ajoelhou na frente dele. Ficaram se olhando. Ela não disse nada. Ele também não. Só balançou a cabeça.

\noindent Ela se levantou. Pegou minha mão.

\noindent \textit{\textquotedblleft{}Anda, Joaquim.\textquotedblright{}}

\noindent Eu olhei para trás. Ele ainda estava sentado. A cabeça caída. O braço estendido no colo.

\noindent Minha mãe puxou minha mão.

\noindent \textit{\textquotedblleft{}Anda.\textquotedblright{}}

\noindent A gente continuou andando.

\noindent A fila na frente. A fila atrás.

\noindent Eu no meio.

\noindent Aprendendo que o mundo é um lugar onde as pessoas ficam para trás e a gente continua andando.

\vspace{0.5em}
\noindent\textit{\textbf{↪ Rotas:} \rota{MEM-03} \rotasep \rota{MEM-01} \rotasep \rota{MEM-06}}